% ===================================================================
\section{4D Topology Programming}
\label{sec:tesseract}

\subsection{Experimental Design}

The tesseract experiment (T-001) uses $n = 8$ channels with pair
specification $\texttt{pairs}_{\text{tess}}(8)$
(Eq.~\ref{eq:tess_pairs}). Two predictions frame the result:

\begin{itemize}
  \item \textbf{Prediction A (4D programming works):} The bridge
    develops four independent $2 \times 2$ blocks, with co-axial
    coupling concentrated within each tesseract axis pair and
    cross-axial coupling suppressed. The eigenvalue spectrum shows a
    $4{+}4$ split: four eigenvalues near zero (suppressed cross-axial
    connections) and four elevated (active co-axial connections).
  \item \textbf{Prediction B (3D attractor):} The bridge collapses to
    three dominant blocks regardless of the 4-pair specification.
    Three-dimensional geometry would be an intrinsic structural limit,
    and the Steersman would be RD-specific rather than general.
\end{itemize}

Training uses TinyLlama-1.1B with identity initialization and
channel size $s = 3$ ($24 / 8$). Four same-seed runs were conducted
(Section~\ref{sec:reproducibility}): T-001r1 (original; crashed at
${\sim}$step~2{,}700 during a March-2026 machine failure, with its
first 1{,}600~steps recovered from its training log), T-001r1b (a
fresh same-seed relaunch, 7{,}100~steps), T-001r2 (complete,
10{,}000~steps), and T-001r3 (a fourth complete run, 10{,}000~steps,
executed on separate hardware).
% Dated edit 2026-07-10: T-001r3 (fourth run) added as the fourth same-seed
% replicate. Source: results/T-001-full-r3/results.json + EXPERIMENT_TRACKER
% 2026-07 sprint update (L-006 dated addition).

\subsection{Results: Prediction A Confirmed}

The Steersman programs 4D tesseract geometry.
Table~\ref{tab:tesseract} reports the primary metrics.

\begin{table}[H]
\centering
\caption{Tesseract contrastive results ($n = 8$). $\rho$ = co-axial
  to cross-axial coupling ratio. Fiedler = Bridge Fiedler eigenvalue
  (mean across adapters). Val loss reported at final step.}
\label{tab:tesseract}
\small
\begin{tabular}{lcccccc}
\toprule
Run & Steps & Fiedler & $\rho$ & Val Loss & Eigenvalue Split & BD \\
\midrule
T-001r1$^{\dagger}$ & 1{,}600 & 0.00055 & 4{,}729:1 & --- & --- & --- \\
T-001r1b & 7{,}100 & 0.000362 & 20{,}944:1 & 0.4049 & $4{+}4$ & Yes \\
T-001r2 & \textbf{10{,}000} & \textbf{0.000191} & \textbf{41{,}564:1} & \textbf{0.4016} & $4{+}4$ & \textbf{Yes} \\
% Dated edit 2026-07-10: T-001r3 fourth run (Hermes RTX 4090); endpoint from
% results/T-001-full-r3/results.json checkpoints[-1] (step 10000).
T-001r3 & 10{,}000 & 0.000193 & 41{,}812:1 & 0.4017 & $4{+}4$ & Yes \\
\bottomrule
\end{tabular}

\smallskip
\raggedright\footnotesize $^{\dagger}$~T-001r1's results file was
overwritten by the T-001r1b relaunch (Section~\ref{sec:reproducibility});
values shown are the last recovered from its training log (step~1{,}600).
Its contemporaneously recorded endpoint (5{,}395:1 at step~2{,}700) has
no surviving artifact and is not tabulated.
\end{table}

Both completed runs produce four independent $2 \times 2$ blocks. The
eigenvalue spectrum exhibits a clean $4{+}4$ split: four eigenvalues
suppressed to $\mathcal{O}(10^{-5})$ (the cross-axial connections)
and four elevated to ${\sim}1.5$--$1.6$ (the co-axial connections within
each block). This is the structural signature of 4D block-diagonal
geometry: each tesseract axis operates as an independent channel
pair, coupled internally and decoupled from the other three axes.

The coupling ratio trajectory rises by three orders of magnitude:
1{,}931:1 at step~400, 5{,}079:1 at step~2{,}800, and \textbf{41{,}564:1}
at step~10{,}000 (T-001r2). Bridge Fiedler at 0.000191 confirms the topological
bifurcation: the spectral attractor at ${\sim}0.09$
(Section~\ref{sec:attractor}) has been broken by the 4D contrastive
specification, collapsing Fiedler by approximately $494\times$
(from $0.0944$ at $n = 8$ spectral-only to $0.000191$). This
collapse is of the same order as the $1{,}020\times$ bifurcation
at $n = 6$ ($0.0918 \to 0.00009$).

Validation loss at 0.4016 is within 0.17\% of the baseline,
consistent with Paper~3's finding that topology programming does not
degrade task performance. The lm-eval benchmark suite confirms this
at the downstream level: the original T-001r1 adapter (2{,}700~steps,
merged and evaluated on 2026-03-16, before the relaunch overwrote its
source artifact; the evaluation outputs survive) shows a mean delta of
$-0.75\%$ across six 0-shot tasks
(HellaSwag, ARC-Easy, PIQA, WinoGrande, OpenBookQA, BoolQ), with
mixed-direction changes (three tasks slightly down, two slightly up,
one negligible). Bridge topology is orthogonal to task capability.

Prediction~A is confirmed: 4D programming works. The Steersman
programs the dimensionality it is told to program.

\subsection{Reproducibility: Four Same-Seed Runs}
\label{sec:reproducibility}

% Dated edit 2026-07-10: run count three -> four; T-001r3 endpoint paragraph
% added below. Sources: results/T-001-full-r3/results.json checkpoints[-1];
% rel dev vs r2's 41,564 per results/EXPERIMENT_TRACKER.md 2026-07 sprint update.
Four runs of the $n = 8$ tesseract contrastive experiment were
conducted with identical hyperparameters and random seed. The
provenance is disclosed in full. T-001r1, the original run, crashed
at ${\sim}$step~2{,}700 during a March-2026 machine failure; when the
experiment was relaunched fresh under the same seed (T-001r1b,
reaching 7{,}100~steps), the relaunch overwrote T-001r1's results
file. T-001r1's first 1{,}600~steps were subsequently recovered from
its surviving training log. T-001r2 ran to completion
(10{,}000~steps). A fourth run, T-001r3, was executed later on
separate hardware (Appendix~\ref{app:details}), also running to
completion (10{,}000~steps) and written to a fresh output directory so
that no existing artifact was overwritten. The first three
trajectories, which carry recoverable per-step feedback logs,
constitute independent same-seed replicates (one partially recovered)
and are cross-correlated below; T-001r3 provides an independent
endpoint confirmation.

\begin{table}[H]
\centering
\caption{Reproducibility of tesseract contrastive training: pairwise
  trajectory agreement between the three same-seed runs with
  recoverable per-step feedback logs, at matched
  feedback steps. T-001r1 pairs span its recovered range
  (steps 200--1{,}600); the T-001r1b~$\times$~T-001r2 pair spans
  their full overlap.}
\label{tab:reproducibility}
\small
\begin{tabular}{lcccc}
\toprule
Pair & $\rho$ Pearson $r$ & $\rho$ max dev & Fiedler $r$ & Steps \\
\midrule
T-001r1 $\times$ T-001r1b & 0.9952 & 10.4\% & 0.9994 & 15 \\
T-001r1 $\times$ T-001r2  & 0.9975 & 5.2\%  & 0.9996 & 15 \\
T-001r1b $\times$ T-001r2 & 0.9984 & 7.2\%  & 0.9999 & 71 \\
\bottomrule
\end{tabular}
\end{table}

Table~\ref{tab:reproducibility} reports the pairwise agreement. All
three pairings track at co-axial/cross-axial ratio correlation
$r \geq 0.995$ and Fiedler correlation $r \geq 0.9994$, with maximum
relative deviations of up to ${\sim}10\%$ at matched steps. The
deviations are the expected signature of floating-point
non-associativity amplified over thousands of optimizer steps; the
correlations show that the topology-programming
\emph{dynamics}---the trajectory shape, the three-order-of-magnitude
ratio growth, the Fiedler collapse---are reproducible across
independent executions. Notably, the run pair that never coexisted
(the recovered T-001r1 and its replacement T-001r1b) agrees nearly as
closely as the pairs that did, confirming that the recovered log
fragment is a faithful record of a genuine same-seed replicate.

T-001r2 continues beyond both earlier runs' termination points,
reaching $\rho = 41{,}564$:1 at step~10{,}000---double T-001r1b's
20{,}944:1 at step~7{,}100---demonstrating continued growth in trend
through the full training run.

% Dated edit 2026-07-10: fourth same-seed run T-001r3 endpoint confirmation.
% Endpoint from results/T-001-full-r3/results.json checkpoints[-1] (co/cross
% 41811.71, Fiedler 1.9347e-4, val 0.4017); rel dev vs r2's 41,564 = 0.596%
% per results/EXPERIMENT_TRACKER.md 2026-07 sprint update.
A fourth same-seed run, T-001r3, was completed on 2026-07-08 on
separate hardware---the Hermes RTX~4090, whereas T-001r2 ran on the
local RTX~6000 Ada (Appendix~\ref{app:details}). Run to the full
10{,}000~steps in a fresh output directory, it reaches
$\rho = 41{,}812$:1 with Bridge Fiedler $1.93 \times 10^{-4}$ at the
final step, reproducing T-001r2's step-10{,}000 endpoint
(41{,}564:1) to within $0.596\%$ relative deviation. That two
independent same-seed executions on \emph{different} GPUs converge on
the coupling ratio to better than one part in 150 is a stronger
reproducibility statement than same-hardware agreement alone: the
programmed 4D block-diagonal endpoint is a property of the pair
specification and the optimizer, not of a particular device's
floating-point environment.

\subsection{Comparison: RD vs.\ Tesseract}

Table~\ref{tab:rd_vs_tess} compares the two confirmed polytope
topologies. Both produce extreme coupling ratios with sub-attractor
Fiedler values and clean eigenvalue splits matching the specified
dimensionality.

\begin{table}[H]
\centering
\caption{Comparison of confirmed polytope topologies. Both use
  contrastive loss with polytope-derived co-axial pair
  specifications. The RD result (H-ch6) is from Paper~3.}
\label{tab:rd_vs_tess}
\small
\begin{tabular}{lcccccc}
\toprule
Geometry & $n$ & Pairs & $\rho$ & Fiedler & Split & Val Loss \\
\midrule
RD (3D) & 6 & 3 co / 12 cross & 70{,}404:1 & 0.00009 & $3{+}3$ & 0.4015 \\
Tesseract (4D) & 8 & 4 co / 24 cross & 41{,}564:1 & 0.000191 & $4{+}4$ & 0.4016 \\
\bottomrule
\end{tabular}
\end{table}

The RD achieves higher $\rho$ (70{,}404 vs.\ 41{,}564) and lower
Fiedler ($0.00009$ vs.\ $0.000191$), likely reflecting the more
favorable co-axial/cross-axial pair ratio at $n = 6$ ($3/12 = 0.25$)
vs.\ $n = 8$ ($4/24 = 0.167$): fewer cross-axial entries to suppress
at $n = 6$ means faster convergence to the block-diagonal attractor.
Both achieve validation loss within 0.03\% of each other and within
0.17\% of the spectral-only baseline.
